\setcounter{equation}{0}
\section{Physical Background}

Quantum mechanics is a branch of physics which models the behavior of energy and matter at atomic and subatomic scales. Its name comes from \emph{quanta}\footnote{the discrete packets of energy proposed by Max Planck in 1900}.

Rather than enumerate informal ``theories'' of quantum mechanics, the cleanest way to present the subject is through its four postulates. Every result we use in this thesis --- superposition, entanglement, the reversibility of computation, measurement statistics --- is a consequence of these four statements and not an additional independent principle.

\subsection{The four postulates}

\paragraph{Postulate 1 (State space).} The state of an isolated quantum system is a unit vector in a complex Hilbert space $\mathcal{H}$. For a single qubit, $\mathcal{H} = \mathbb{C}^2$, and the state is
\begin{equation}
\ket{\psi} = \alpha\ket{0} + \beta\ket{1}, \qquad |\alpha|^2 + |\beta|^2 = 1.
\end{equation}
Two states that differ only by a global phase $e^{i\theta}$ are physically indistinguishable.

\paragraph{Postulate 2 (Evolution).} The evolution of a closed quantum system is unitary. A discrete time step is described by a unitary operator $U$ acting on the state:
\begin{equation}
\ket{\psi'} = U\ket{\psi}.
\end{equation}
Continuous evolution is governed by the Schr\"{o}dinger equation $i\hbar\,\partial_t\ket{\psi} = H\ket{\psi}$, whose finite-time solution is $U = e^{-iHt/\hbar}$. For computation we usually work with $U$ directly and treat the underlying Hamiltonian $H$ as an implementation detail of the physical hardware.

\paragraph{Postulate 3 (Measurement).} A measurement in the computational basis returns outcome $x$ with probability $|\braket{x}{\psi}|^2$ and projects the state onto $\ket{x}$. More generally, a measurement is described by a collection of operators $\{M_m\}$ satisfying $\sum_m M_m^\dagger M_m = I$; outcome $m$ occurs with probability $\bra{\psi} M_m^\dagger M_m \ket{\psi}$, and after observing $m$ the state collapses to $M_m \ket{\psi}/\sqrt{\bra{\psi} M_m^\dagger M_m \ket{\psi}}$.

\paragraph{Postulate 4 (Composite systems).} The state space of a composite system is the tensor product of the component state spaces. For two qubits,
\begin{equation}
\mathcal{H} = \mathbb{C}^2 \otimes \mathbb{C}^2 = \mathbb{C}^4.
\end{equation}

\subsection{Superposition is linearity}

Superposition is the statement that if $\ket{\psi_1}$ and $\ket{\psi_2}$ are valid states, then so is any normalised linear combination $\alpha\ket{\psi_1} + \beta\ket{\psi_2}$. It is a direct consequence of Postulate 1: state space is a vector space.

This has nothing to do with the Heisenberg uncertainty principle. The uncertainty principle is a statement about \emph{observables}: for any two observables $A$ and $B$,
\begin{equation}
\Delta A \cdot \Delta B \;\geq\; \tfrac{1}{2}\,\bigl|\bra{\psi}[A,B]\ket{\psi}\bigr|.
\end{equation}
When $A$ and $B$ do not commute, no quantum state can simultaneously have sharp values for both. Position and momentum are the canonical example. Superposition, in contrast, exists in finite-dimensional Hilbert spaces with no position observable at all (a single qubit being the simplest example). The two ideas live on different shelves and should not be conflated.

\subsection{Entanglement and coherence are distinct}

A two-qubit state is \emph{entangled} if and only if it cannot be written as a tensor product $\ket{\phi}\otimes\ket{\chi}$ of single-qubit states. The four canonical Bell states are all maximally entangled:
\begin{equation}
\ket{\Phi^\pm} = \tfrac{1}{\sqrt{2}}\bigl(\ket{00} \pm \ket{11}\bigr), \qquad
\ket{\Psi^\pm} = \tfrac{1}{\sqrt{2}}\bigl(\ket{01} \pm \ket{10}\bigr).
\end{equation}
Measuring the first qubit of $\ket{\Phi^+}$ yields outcomes that \emph{match} the second qubit's; measuring the first qubit of $\ket{\Psi^-}$ yields outcomes that anti-match. Entanglement is non-classical correlation in general --- not specifically anti-correlation, and not specifically the spin-singlet case.

\emph{Coherence} is a different concept entirely. A state is coherent (with respect to a chosen basis) when its density matrix has non-zero off-diagonal elements in that basis. For a single qubit in the state $\ket{+} = (\ket{0}+\ket{1})/\sqrt{2}$, the density matrix
\begin{equation}
\rho_+ = \tfrac{1}{2}\begin{pmatrix} 1 & 1 \\ 1 & 1 \end{pmatrix}
\end{equation}
is coherent in the $\{\ket{0},\ket{1}\}$ basis --- the off-diagonal elements record the relative phase between $\ket{0}$ and $\ket{1}$. After decoherence in that basis, the diagonal survives but the off-diagonals vanish:
\begin{equation}
\rho_{\text{dec}} = \tfrac{1}{2}\begin{pmatrix} 1 & 0 \\ 0 & 1 \end{pmatrix}.
\end{equation}
A single-qubit superposition is coherent but trivially unentangled; one can also have entanglement on noisy mixed states with reduced coherence. They are separate resources. Decoherence --- caused by uncontrolled interaction with the environment, sometimes by even a single stray photon --- is the process that destroys both.

\subsection{Why computation must be reversible}

Landauer (1961) \cite{landauer1961} showed that erasing one bit of information in a system at temperature $T$ dissipates at least $k_B T \ln 2$ joules of heat. Bennett (1973) \cite{bennett1973} showed that any classical computation can be made logically reversible by retaining intermediate results, so heat dissipation is not fundamentally required for computation.

Quantum computation is reversible by construction: Postulate 2 says that closed-system evolution is unitary, and unitary maps are invertible ($U^{-1} = U^\dagger$). The deeper reason this matters, beyond thermodynamic prudence, is that non-unitary evolution is what decoherence \emph{looks like} from the outside. A computation that ``leaks'' information into the environment is, by definition, a computation entangled with environmental degrees of freedom that we cannot track, and tracing them out collapses our superposition. Reversibility is therefore the property a quantum computer must preserve to keep its computational state coherent.

In practice, achieving reversibility requires extensive isolation of the qubit register, and active mitigation through quantum error-correcting codes, which encode each logical qubit redundantly across many physical qubits.

\subsection{The no-cloning theorem}

One result deserves explicit mention because it shapes both quantum algorithms and quantum error correction: there is no unitary $U$ such that
\begin{equation}
U\ket{\psi}\ket{0} = \ket{\psi}\ket{\psi}
\end{equation}
for arbitrary $\ket{\psi}$. The proof is one line: assume such a $U$ exists, apply it to two non-orthogonal states $\ket{\psi}$ and $\ket{\phi}$, take inner products on both sides of the resulting identity, and conclude $\braket{\psi}{\phi} = \braket{\psi}{\phi}^2$, which forces $\braket{\psi}{\phi} \in \{0,1\}$ --- a contradiction \cite{wootterszurek1982}. This is why quantum error correction cannot simply ``copy and check'' to recover from errors; it must protect the state without ever reading it directly.

In order to explain quantum algorithms and the efficiency gains they achieve, we now need to understand how the abstract qubit register works internally, and how it is described mathematically.
